\documentclass[10pt,a4paper]{amsart}
\usepackage[margin=19mm]{geometry}
\usepackage{amsmath,amssymb,mathtools}
\usepackage[scr=rsfs,cal=euler]{mathalfa}
\usepackage{xfrac}
\usepackage[unicode,hidelinks,bookmarks=false]{hyperref}
\newcommand{\ie}{i.e.\ }
\newcommand{\cf}{cf.\ }
\newcommand{\resp}{respectively\ }
\newcommand{\Subsection}{\subsection}
\providecommand{\nobreakdash}{\nobreak\hbox{-}}
\setlength{\emergencystretch}{2em}
\multlinegap0pt
\usepackage{mdframed}
\usepackage{mdframed}
\usepackage{mdframed}
\usepackage{mdframed}
\usepackage{amssymb}
\usepackage{mathtools}
\usepackage{bm}
\usepackage{xcolor}
\usepackage{float}
\usepackage{algorithm}\usepackage{algpseudocode}
\usepackage{mdframed}
\usepackage[shortlabels]{enumitem}
\usepackage{tikz}
\newtheorem{theorem}{Theorem}
\title{On the Extension Theorem for Packing Steiner Forests}
\author{Jinghan A Zeng}
\date{Source: arXiv:2603.16956v2 (CC BY 4.0)}
\begin{document}
\maketitle

\noindent\emph{Source and licence.} On the Extension Theorem for Packing Steiner Forests. Version arXiv:2603.16956v2 (CC BY 4.0), distributed by arXiv under the Creative Commons Attribution 4.0 International licence, http://creativecommons.org/licenses/by/4.0/, as recorded in the arXiv licence metadata for this exact version and shown on the abstract page https://arxiv.org/abs/2603.16956v2. Full licence notice: \texttt{licenses/arxiv-2603.16956-CC-BY-4.0.txt}.\par\smallskip

\noindent\textit{Editorial scope.} Counted unit: the article's theorem theorem (source0.tex lines 776--778). The article's complete original proof of it is delivered with it (source0.tex lines 780--796, one \texttt{proof} environment of 17 lines). No mathematics is written, repaired or re-derived: the statement and the proof are the article's own lines, in the article's own notation, and no retained line is substituted or re-flowed.  No further source-local context is needed: every cross-reference of every retained line resolves inside the two delivered spans.  Nothing was omitted from any retained span, so the package carries no omission record.  No retained line contains a citation, so the wrapper carries no bibliography and no bibliography entry.  No retained line contains a figure environment, an image-inclusion command or drawing code, so no image is delivered and the package declares no asset dependency.  Editorial formatting only, in addition: an AMS \texttt{amsart} preamble that restates the article's own declarations and macros used by the retained lines, namely the environments \texttt{theorem} on separate counters, together with the standard packages those lines need.  The retained environments print in this document as Theorem 1 in order of appearance, whereas the article prints the same items with the numbering of its own full document.  The complete native source of the article as distributed in the arXiv e-print for this exact version is bundled unchanged as \texttt{sources/196586/source0.tex}.
\begin{theorem}[Steiner Forest Packing]
    Given a loopless multigraph $G$, $S_1,S_2,...,S_t, R \subseteq V(G)$ where $(S_1 \cup S_2...S_t) \cap R = \emptyset$, such that each $S_i$ is $32k$-edge-connected, $|S_i| \geq 2$, and all vertices in $R$ are incident to at least $30k$ edges. Then there exist $k$ edge-disjoint $(S_1,S_2,...,S_t)$-subgraphs that balance $R$. 
\end{theorem}

\begin{proof}
    Suppose, for contradiction, that this is false. Then take a counterexample with $|V(G)|+|E(G)|$ minimized. We can assume that $S_1 \cup S_2...S_t$ is not $32k$-edge-connected, otherwise this problem reduces down to the previous lemma.

    Let $C$ be a subset of $V(G)$ such that the number of edges between $C$ and $V(G)-C$ is at most $32k$, for all $1 \leq i \leq t$, either $S_i \subseteq C$ or $S_i \subseteq V(G)-C$, and there exists $i,j$ such that $S_i \subseteq C$ and $S_j \subseteq V(G)-C$. We know such a $C$ must exist, since there is an $S_1 \cup S_2...S_t$-cut of less than $32k$, which keeps all $S_i$ intact since each are $32k$-edge connected. Pick a $C$ with the minimum possible number of vertices, and we will use $C_2$ to denote $V(G)-C$. Contract $C_2$ into a vertex $v_2$ and $C$ into a vertex $v$ to obtain $G_2 = C_2+v$ and $G_1 = C+v_2$. 

    If $v$ has degree at least $30k$, add it to $R$ so we can ensure its balancing. Note that since $|C| \geq 2$, $G_2$ is strictly smaller than $G$, so we can inductively find $k$ edge disjoint subgraphs, each of them which connect each $S_i \subseteq C_2$ and balance each $R \cap C_2$ on $G_2$. 

    Let $S^*$ denote $(S_1 \cup S_2...S_t) \cap C$. We first claim $S^*$ is $32k$-edge-connected in $G_1$. Indeed, if there is a cut of $< 32k$ that splits $C$ into $C_a$ and $C_b$, it must keep all $S_i \subseteq C$ intact. Without loss of generality suppose $v_2 \in C_b$, then we can replace $C$ with $C_a$ to get a smaller set that satisfies the desired properties of $C$, contradicting the minimality of $C$.

    We next claim that there are $d(v_2)$ edge disjoint paths from $v_2$ to any $s \in S^*$. Suppose, for contradiction, that there was a cut of $< d(v_2)$ with $s \in C_a$ and $v_2 \in C_b$. Then clearly, $S^* \subseteq C_a$, and $v_2$ cannot be a single vertex since its degree is $d(v_2)$, so $|C_b| \geq 2$. Then this again contradicts the minimality of $C$. Based on the forest packing on $G_2$, copy the corresponding labels over to $v_2$ on $G_1$. Then, pick any vertex $s' \in S^*$, and add $32k-d(v_2)$ fake edges from $v_2$ to $s'$. 

    We claim that $v_2$ is now $32k$-edge-connected to any $s \in S^*$. Indeed, suppose there is a cut of $< 32k$ separating $v_2$ and $s$. Since the cut is less than $32k$, $s'$ and $s$ will still be connected, so it means that all the fake edges must have been deleted. Then this implies that the deletion of less than $d(v_2)$ edges disconnects $v_2$ and $s$, contradicting what we found previously. Assign the fake edges to have no label. Note that if $d(v_2) \geq 30k$, then $v_2$ is balanced as $v$ is balanced with respect to the packing on $G_2$. Otherwise, $v_2$ is balanced with the addition of the fake edges, since if $d(v_2) < 30k$, we added at least $2k$ fake edges. 

    Apply the extension theorem to get $k$-edge-disjoint $S^* \cup v_2$ subgraphs $H_1,H_2,...,H_k$ that balances $R \cap C$. We now delete the fake edges we added from the graph entirely. Furthermore, to make this extension a strict edge extension, if an edge incident to $v_2$ was placed in $H_i$ in the packing, yet it was not initially assigned a label of $i$, we remove the edge from $H_i$ but keep it in the graph. Note that since each subgraph $H_i$ is still connected after the deletion of $v_2$, all of $H_i-v_2$ is still connected after the deletion or removal of those edges. Furthermore, since none of the fake edges were incident to a vertex in $R$, and removing edge labels does change a balanced vertex to a unbalanced vertex, $C \cap R$ is still balanced. Let $H_1',H_2',...,H_k'$ be the $k$ edge disjoint subgraphs we found on $G_2$. We note that $H_1 \uplus H_1', H_2 \uplus H_2'... H_k \uplus H_k'$ balances $R$ and are edge disjoint by construction. We claim that $H_i \uplus H_i'$ is a $(S_1,S_2,...,S_t)$-subgraph. 

    Note that if any $S_j \subseteq C$, then $S_j$ is connected in $H_i \uplus H_i'$ as $H_i-v_2$ is connected. For any $S_j \subseteq C_2$, given $a, b \in S_j$, if there exists a path from $a$ to $b$ that does not traverse through $v$ in $H_i'$, then $a$ and $b$ are connected on $H_i'$ as well. Otherwise, if all paths from $a$ to $b$ traverse through $v$ in $H_i'$, take a path $P_1$ from $a$ to $v$ on $H_i'$, and a path $P_2$ from $v$ to $b$ on $H_i'$. Let $e$ be the last edge of $P_1$ and $e'$ be the first edge of $P_2$. Then by construction, $e$ and $e'$ appear on $H_i$ (even after deletion of the fake edges and removal of edges to get a strict extension). Setting $e = uv_2$ and $e' = u'v_2$, since $H_i-v_2$ is connected, we can find a path from $e$ to $e'$ on $H_i-v_2$ and hence $H_i \uplus H_i'$. Then, we can traverse from $a$ to $b$ on $H_i \uplus H_i'$ by first following the $P_1$ from $a$ to $u$, traversing from $u$ to $u'$, and then following $P_2$ to reach $b$. Hence, we have found $k$ edge-disjoint $(S_1,S_2,...,S_t)$-subgraphs that balance $R$, as desired. 
\end{proof}

\end{document}
